\section{Discussion}
\label{sec:discussion}

{\bf Why does more detail not hurt 7B and 14B models?}
Our longest prompts are 1{,}000--1{,}700 tokens, far below the multi-thousand-token regimes where context-length degradation has been reported~\citep{liu2024lost,levy2024same}.
Our \exhaustive tiers are also consistent: they elaborate the task and never contradict it.
Under these conditions, extra rules are either redundant (classification, reasoning) or mildly useful (span boundaries in extraction, ``lead with the most newsworthy fact'' in summarization).
The main lever is the first step from \minimal to a short, specific instruction; beyond that, length is close to neutral.
The cost of a long instruction shows up elsewhere: the 7B model doubles its reasoning length when told to verify its work, which raises latency and cost without changing accuracy.

{\bf What goes wrong in tiny models?}
The 0.5B model's failure is a shift in its default label, not a loss of the format.
Two observations point to the text immediately before the input as the cause.
First, \padded has exactly the same task content as \standard but recovers most of the accuracy; the only difference is that the article follows neutral prose instead of a rule list ending in ``Output exactly one of \ldots (World, Sports, Business, Sci/Tech).''
Second, moving the instruction to the system message removes most of the drop at \standard.
For a 0.5B model, rules about companies, markets, and technology placed right before the article seem to prime the ``Business'' label, much like the recency bias reported for few-shot demonstrations~\citep{zhao2021calibrate}.
This is an interference effect of content and position, not a dilution effect of length.
A residual decline of about 12 points under system placement shows that small models are also harmed by detailed rule content itself.

{\bf Practical implications.}
Our results support four simple rules.
(1) Always include a short, task-specific definition; \minimal prompts leave about 10 F1 points on the table for extraction.
(2) For $\geq$3B instruction-tuned models, extra detail is cheap insurance with no measurable accuracy cost, but it can lengthen outputs.
(3) Restate the output format at the end of a long prompt; this made \exhaustive the most compliant tier for reasoning.
(4) For sub-2B models, prefer minimal instructions, or put them in the system message.

\subsection{Limitations}
\label{sec:limitations}

\para{One model family.}
All five models are \qwen. The planned Llama-3.1-8B and \qwen-32B runs could not be completed because the shared GPUs were occupied, so conclusions about ``larger models'' rest on 14B alone and cross-family generality is untested.

\para{One wording per tier.}
Within a tier, wording and length are confounded.
The nested design and the padding control reduce this problem but do not remove it.
A different \exhaustive prompt, for example one with conflicting or over-specific rules, could well hurt; ours were written to be realistic and consistent.

\para{Imperfect padding control.}
\padded is word-matched, which gives 7--12\% fewer tokens than \exhaustive.
It also changes what precedes the input, so the \padded\versus\standard comparison mixes length with recency.

\para{Statistical power.}
With $n=252$ per cell, binary-metric CIs are about $\pm$5--6 points, so the power to detect a 3-point effect in a single cell is low.
The size correlation uses only five points.

\para{Metrics and decoding.}
ROUGE-L is a coarse proxy for summary quality, and we did not run human or LLM-judge evaluation.
\agnews labels are noisy, which caps accuracy at roughly 85--90\%.
Compliance definitions are strict and chosen by us.
We use greedy decoding with a single seed and capture variability only through item resampling.
All four datasets are public and likely seen in pre-training, which affects absolute levels more than within-model comparisons across tiers.

\para{Broader impact.}
This work aims to help practitioners write better prompts and does not introduce new capabilities.
A risk is over-generalization: our finding that detail is harmless for $\geq$3B models comes from one family at moderate prompt lengths and should not be read as a guarantee for every model or for very long agent prompts.
